\documentclass[10pt,a4paper]{amsart}
\usepackage[margin=19mm]{geometry}
\usepackage{amsmath,amssymb,mathtools}
\usepackage[scr=rsfs,cal=euler]{mathalfa}
\usepackage{xfrac}
\usepackage[unicode,hidelinks,bookmarks=false]{hyperref}
\newcommand{\ie}{i.e.\ }
\newcommand{\cf}{cf.\ }
\newcommand{\resp}{respectively\ }
\newcommand{\Subsection}{\subsection}
\providecommand{\nobreakdash}{\nobreak\hbox{-}}
\setlength{\emergencystretch}{2em}
\multlinegap0pt
\usepackage{amssymb}
\usepackage[shortlabels]{enumitem}
\usepackage{tikz}
\newtheorem{lemma}{Lemma}
\usepackage{cleveref}
\crefname{lemma}{Lemma}{Lemmas}
\title{Box-Delaunay graphs of large chromatic number}
\author{Istv\'an Tomon}
\date{Source: arXiv:2609.01439v1 (CC BY 4.0)}
\begin{document}
\maketitle

\noindent\emph{Source and licence.} Box-Delaunay graphs of large chromatic number. Version arXiv:2609.01439v1 (CC BY 4.0), distributed by arXiv under the Creative Commons Attribution 4.0 International licence, http://creativecommons.org/licenses/by/4.0/, as recorded in the arXiv licence metadata for this exact version and shown on the abstract page https://arxiv.org/abs/2609.01439v1. Full licence notice: \texttt{licenses/arxiv-2609.01439-CC-BY-4.0.txt}.\par\smallskip

\noindent\textit{Editorial scope.} Counted unit: the article's lemma lemma (source0.tex lines 238--241). The article's complete original proof of it is delivered with it (source0.tex lines 243--260, one \texttt{proof} environment of 18 lines). No mathematics is written, repaired or re-derived: the statement and the proof are the article's own lines, in the article's own notation, and no retained line is substituted or re-flowed.  Retained uncounted as the source-local context the proof needs: lines 140--142.  Nothing was omitted from any retained span, so the package carries no omission record.  No retained line contains a citation, so the wrapper carries no bibliography and no bibliography entry.  No retained line contains a figure environment, an image-inclusion command or drawing code, so no image is delivered and the package declares no asset dependency.  Editorial formatting only, in addition: an AMS \texttt{amsart} preamble that restates the article's own declarations and macros used by the retained lines, namely the environments \texttt{lemma} on separate counters, together with the standard packages those lines need.  The retained environments print in this document as Lemma 1 in order of appearance, whereas the article prints the same items with the numbering of its own full document.  The complete native source of the article as distributed in the arXiv e-print for this exact version is bundled unchanged as \texttt{sources/209110/source0.tex}.
\begin{lemma}\label{lemma:gadget}
For $I\subset [m]$, the probability that $I$ is independent in the Hasse diagram of $\alpha$ is  $(|I|+1)2^{-|I|}$.
\end{lemma}

\begin{lemma}
Let $I\subset [n]$. The probability that $I$ is an independent set in $G$ is at most 
$$2^{-|I|k}\left(\frac{|I|m}{n}+1\right)^{kn/m}.$$
\end{lemma}

\begin{proof}
We want to argue that the events $\{I\cap K\text{ is an independent set in the Hasse diagram of }\pi|_K\}$ are independent for the different slices $K$. As the slices themselves are dependent, we need to be somewhat careful with this statement, so we present a formal argument. 

By the previous, if $I$ is an independent set, then for every slice $K$, $I\cap K$ is an independent set in the Hasse diagram of $\pi|_{K}$. We denote by $H_K$ the Hasse diagram of $\pi|_{K}$. For $j=1,\dots,k$, let $\mathcal{A}_j$ be the event that $I\cap K$ is an independent set in $H_K$ for every level-$j$ slice $K$. Then
$$\mathbb{P}(I\text{ is an independent set in }G)\leq \mathbb{P}(\mathcal{A}_1\cap\dots\cap\mathcal{A}_k)=\prod_{j=1}^{k}\mathbb{P}(\mathcal{A}_j|\mathcal{A}_1,\dots,\mathcal{A}_{j-1}).$$
Here, the events $\mathcal{A}_1,\dots,\mathcal{A}_{j-1}$ are fully determined by $\pi_{j-1}$. Fixing any outcome $\tau$ of $\pi_{j-1}$, the level-$j$ slices are also fully determined, let $L_{a,s}^{(j)}$ denote the corresponding outcome of the slice $K_{a,s}^{(j)}$. Then we have
\begin{align*}
    \mathbb{P}(\mathcal{A}_j|\pi_{j-1}=\tau)&=\prod_{L} \mathbb{P}(I\cap L\text{ is an independent set in }H_{L})\\
    &=\prod_{L}  2^{-|I\cap L|}(|I\cap L|+1)\\
    &=2^{-|I|}\prod_{L}(|I\cap L|+1),
\end{align*}
where the product ranges over all level-$j$ slices $L=L_{a,s}^{(j)}$. The first equality holds by the independence of the permutations $\beta^{(j)}_{a,s}$, the second equality holds by \Cref{lemma:gadget}, and the third equality holds by noting that the level-$j$ slices partition $[n]$. We further bound the right-hand-side by the AM-GM inequality. As the number of level-$j$ slices is $n/m$, we can write
$$\prod_{L}(|I\cap L|+1)\leq \left(\frac{\sum_{L} |I\cap L|+1}{n/m}\right)^{\frac{n}{m}}=\left(\frac{|I|m}{n}+1\right)^{\frac{n}{m}}.$$
Thus, filtering the event $\mathcal{A}_1\cap\dots\cap\mathcal{A}_{j-1}$ by the possible outcomes $\tau$ of $\pi_{j-1}$, we get
$$\mathbb{P}(\mathcal{A}_j|\mathcal{A}_1,\dots,\mathcal{A}_{j-1})\leq 2^{-|I|} \left(\frac{|I|m}{n}+1\right)^{\frac{n}{m}}.$$
In conclusion,
$$\mathbb{P}(I\text{ is an independent set in }G)\leq\prod_{j=1}^{k}\mathbb{P}(\mathcal{A}_j|\mathcal{A}_1,\dots,\mathcal{A}_{j-1})\leq 2^{-k|I|} \left(\frac{|I|m}{n}+1\right)^{\frac{kn}{m}}.$$
\end{proof}

\end{document}
